\scanpage{32}
\mspage{6}
il n'est pas possible d'y introduire une métrique riemannienne canonique, même modulo cte. Considérons cependant la donnée d'une structure supplémentaire du type
\[
H\subset\operatorname{Aut}_{\mathrm{conf}}(X)=G
\]
où $H$ est un sous-groupe \uline{compact} de $\operatorname{Aut}_{\mathrm{conf}}(X)$ (condition \uline{nécessaire} pour que $H$, supposé fermé, invarie une métrique euclidienne sphérique comp. avec la str. conforme). Soit $H^+=H\cap\operatorname{Aut}^+(X)$ (groupe des autom. qui conservent l'orientation). OPS $X=\mathbb P(V)$, $V$ vectoriel complexe de rg $2$, \deleted{\unsure{étant}} soit $\unsure{\widetilde H|\widehat H}$ l'image inverse de $H^+$ dans $\unsure{\operatorname{SL}(V,V)|\operatorname{SL}(V)}$, c'est donc un sous-groupe cpt, qui invarie nécessairement une forme hermitienne déf. pos. sur $V$, d'où la métrique sphérique ass. sur $X$.\inserted{(normalisée par la condition $\int_X dx=1$)}, Dans le cas où \deleted{cette métrique}\inserted{la métrique} hermitienne invariante est unique mod. cte, on trouve donc une métrique\inserted{sphérique normalisée} canonique sur $X$, définie par la\inserted{seule} donnée de $H^+$. Cette métrique est $=$ invariante par $\operatorname{Norm}_G(H^+)$\deleted{\unsure{de $H$}}, et a fortiori par $H$, et à nouveau caractérisée par cette condition.
% archive foliation: 34
